% !TeX root = main.tex
% !TEX program = pdflatex

\chapter{Metrología y Caracterización de Sistemas de Medida}

\section{El Sistema de Medida de Ingeniería}\label{sec:sistema-medida}

Un sistema de instrumentación no es solo un sensor; es una cadena de bloques que transforma una variable física en una información útil~\cite{WHO2025}. Esta cadena incluye:

\begin{itemize}
  \item Mundo Físico: La magnitud a medir (temperatura, presión, flujo, etc.).
  \item Sensor/Transductor: Convierte la magnitud física en eléctrica (ej. Resistencia, Voltaje).
  \item Acondicionamiento de Señal: Filtrado, amplificación, linealización.
  \item Adquisición de Datos (DAQ): Conversión Analógica-Digital (ADC).
  \item Procesamiento y Presentación: Microcontrolador/PC y visualización. 
\end{itemize}

Incidir en que el eslabón más débil de la cadena determina la precisión de todo el sistema. Un sistema de medida es tan bueno como su componente más impreciso.

Un sistema de instrumentación moderno es una cadena de transferencia de energía e información. Cada etapa introduce una degradación de la señal (ruido) y una incertidumbre asociada.
Bloques Funcionales
\begin{enumerate}
    \item Transductor/Sensor: Interfaz con el mundo físico. Convierte la magnitud física $x(t)$ (temperatura, presión) en una señal eléctrica $s(t)$
    \item Acondicionador de Señal: Mejora la calidad de la señal (amplificación, filtrado, linealización, aislamiento). Su objetivo es maximizar la relación señal-ruido (SNR).
    \item Sistema de Adquisición (DAQ): Compuesto por el S&H (Sample and Hold) y el ADC. Realiza la discretización en tiempo y amplitud.
    \item Procesado y Control: Algoritmos de corrección, promediado y protocolos de comunicación (USB, Ethernet, CAN).
\end{enumerate}

% \begin{figure}[ht]
%     \centering
%     \begin{tikzpicture}[node distance=2.5cm, auto, >=stealth]
%       \node [block] (phys) {Magnitud\\Física $x(t)$};
%       \node [block, right of=phys] (sens) {Sensor /\\Transductor};
%       \node [block, right of=sens] (cond) {Acondic.\\de Señal};
%       \node [block, right of=cond] (daq) {DAQ\\(ADC)};
%       \node [block, right of=daq] (proc) {Procesado /\\PC};
      
%       \draw [->] (phys) -- (sens);
%       \draw [->] (sens) -- node[above] {$\Delta V, \Delta R$} (cond);
%       \draw [->] (cond) -- node[above] {$0-5V$} (daq);
%       \draw [->] (daq) -- node[above] {Digital} (proc);
      
%       \node[below of=cond, node distance=1.2cm] (noise) {Ruido / Interferencias};
%       \draw [->, red, dashed] (noise) -- (cond);

%     \end{tikzpicture}
%     \caption{Cadena de Instrumentación: Desde la magnitud física hasta la información útil.}
%     \label{fig:cadena_instrumentacion}
% \end{figure}
% \begin{figure}[ht]
%     \centering
%     \begin{tikzpicture}[node distance=2.5cm, auto, >=stealth, every node/.style={inner sep=0pt}]
%       \node [block] (phys) {Magnitud\\Física $x(t)$};
%       \node [block, right of=phys] (sens) {Sensor /\\Transductor};
%       \node [block, right of=sens] (cond) {Acondic.\\de Señal};
%       \node [block, right of=cond] (daq) {DAQ\\(ADC)};
%       \node [block, right of=daq] (proc) {Procesado /\\PC};
      
%       \draw [->] (phys) -- (sens);
%       \draw [->] (sens) -- node[above] {$\Delta V, \Delta R$} (cond);
%       \draw [->] (cond) -- node[above] {$0-5V$} (daq);
%       \draw [->] (daq) -- node[above] {Digital} (proc);
      
%       \node[below of=cond, node distance=1.2cm] (noise) {Ruido / Interferencias};
%       \draw [->, red, dashed] (noise) -- (cond);
      
%     \end{tikzpicture}
%     \caption{Cadena de Instrumentación: Desde la magnitud física hasta la información útil.}
%     \label{fig:cadena_instrumentacion}
% \end{figure}

\section{Terminología según el VIM (Vocabulario Internacional de Metrología)}
\label{sec:terminologia-vim}

Es vital que se utilice la terminología profesional:

\begin{itemize}
  \item \textbf{Magnitud}: Propiedad física que se puede medir (ej. Temperatura, Longitud).
  \item \textbf{Medida}: Proceso de asignar un valor numérico a una magnitud.
  \item \textbf{Valor Verdadero}: Valor ideal sin errores (inaccesible en la práctica).
  \item \textbf{Valor Medido}: Resultado obtenido del proceso de medida.
  \item \textbf{Error de Medida}: Diferencia entre el valor medido y el valor verdadero o de referencia. $e = V_{\text{med}} - V_{\text{ref}}$
  \item \textbf{Exactitud (Accuracy)}: Grado de concordancia entre el valor medido y el valor verdadero.
  \item \textbf{Precisión (Precision)}: Grado de concordancia entre resultados repetidos.
  \item \textbf{Resolución}: Mínimo cambio detectable por el sistema de medida. Ejemplo: En un ADC de \SI{10}{bits} con rango \SI{0}{V} a \SI{5}{V}, la resolución es \SI{5}{V}/$2^{10}$ bits  \(\approx\) \SI{0.05}{V/\bit}.
  \item \textbf{Sensibilidad}: Relación entre el cambio en la salida y el cambio en la magnitud medida.Pendiente de la curva de transferencia ($S = \frac{\Delta \text{Output}}{\Delta \text{Input}}$).
  \item \textbf{Linealidad}: Grado en que la salida es proporcional a la entrada.
  \item \textbf{Repetibilidad}: Capacidad de obtener el mismo resultado bajo las mismas condiciones.
  \item \textbf{Reproducibilidad}: Capacidad de obtener resultados similares bajo condiciones diferentes (operadores, equipos, etc.).
  \item \textbf{Calibración}: Proceso de ajustar el sistema de medida para minimizar errores, comparándolo con un estándar de referencia.
  \item \textbf{Trazabilidad}: Capacidad de relacionar las mediciones con estándares internacionales a través de una cadena documentada de calibraciones.
  \item \textbf{Incertidumbre de Medida}: Rango dentro del cual se espera que se encuentre el valor verdadero con un nivel de confianza específico. Se expresa como $U = k \cdot u_c$, donde $u_c$ es la incertidumbre combinada y $k$ es el factor de cobertura (ej. $k=2$ para aproximadamente 95\% de confianza).
\end{itemize}


